\section{推理算法：logits、softmax、采样、搜索与投机解码}

\subsection{模型输出的不是文字，而是分数分布}

decoder 最后一层输出 hidden 后，\code{lm\_head} 会得到词表大小的 logits。logits 是每个 token 的未归一化分数，不是概率，也不是最终文字。推理算法的任务是把 logits 变成下一个 token id。

\begin{keyidea}
“模型会不会回答得好”不只由权重决定，也由解码算法决定。贪心、温度、top-k、top-p、重复惩罚、beam search 和投机解码都会改变输出、延迟、可复现性和服务成本。
\end{keyidea}

如果把模型主体看成 \code{hidden -> logits}，采样器就是 \code{logits -> next token}。它在数学上很小，在业务上很重要：同一组 logits，贪心可能稳定但死板，高温采样可能多样但不稳定，top-p 可能减少低质量尾部 token，重复惩罚可能降低复读，但也可能破坏术语一致性。

\subsection{softmax 和数值稳定}

softmax 把 logits 转成概率：

\[
p_i = \frac{e^{z_i}}{\sum_j e^{z_j}}
\]

直接计算容易溢出。稳定做法是减去最大值：

\[
p_i = \frac{e^{z_i-z_{max}}}{\sum_j e^{z_j-z_{max}}}
\]

这不会改变概率，因为分子分母同时乘了同一个常数。采样器必须使用稳定 softmax，尤其在低温或 logits 幅度较大时。

\begin{lstlisting}[numbers=none,caption={稳定 softmax 伪代码}]
max_logit = max(logits)
sum = 0
for i in vocab:
  prob[i] = exp(logits[i] - max_logit)
  sum += prob[i]
for i in vocab:
  prob[i] /= sum
\end{lstlisting}

在工程实现中，完整 softmax 会扫描整个词表。若词表有三万到十几万 token，sampler 的成本并非完全可以忽略。GPU 后端若每步把完整 logits 拷回 CPU 再采样，会引入同步和传输成本；CPU 后端则要注意 top-k 选择和概率归一化的复杂度。

\subsection{贪心解码}

贪心解码直接选择最大 logit：

\[
next = \arg\max_i z_i
\]

它的优点是确定、简单、可复现，适合测试、oracle 和某些结构化任务。缺点是缺少多样性，容易陷入重复，且不一定得到全局最优序列。贪心解码不需要 softmax，因为最大 logit 对应最大概率。

\begin{lstlisting}[numbers=none,caption={贪心解码}]
best_id = 0
best_value = logits[0]
for i in 1..vocab:
  if logits[i] > best_value:
    best_value = logits[i]
    best_id = i
return best_id
\end{lstlisting}

reference 和回归测试常用贪心，因为它去掉了随机性。若贪心下 logits top-1 都不一致，先不要讨论采样策略，应该回到模型计算、量化或后端 oracle。

\subsection{温度：控制分布尖锐程度}

温度采样先把 logits 除以温度 \code{T}：

\[
p_i = \operatorname{softmax}(z_i / T)
\]

\code{T} 越小，分布越尖锐，模型更接近贪心；\code{T} 越大，分布越平，低分 token 更容易被采到。温度不能简单理解成“创造力”。它改变的是概率分布熵，过高会增加胡言乱语，过低会增加模板化和重复。

\begin{center}
\begin{tabular}{p{0.18\linewidth}p{0.36\linewidth}p{0.30\linewidth}}
\toprule
温度 & 效果 & 常见用途 \\
\midrule
接近 0 & 接近贪心，稳定 & 测试、代码、结构化输出 \\
中等 & 保留多样性 & 问答、摘要、普通对话 \\
较高 & 分布更平，随机性更强 & 创意任务，但风险高 \\
\bottomrule
\end{tabular}
\end{center}

实现上要防止 \code{T=0} 直接除零。工程中通常把 \code{T<=0} 解释成贪心，或者拒绝参数。

\subsection{top-k：只在前 k 个候选中采样}

top-k 先保留 logit 最大的 \code{k} 个 token，把其他 token 概率置零，再归一化采样。它减少低质量尾部 token 的机会，也让采样成本从全词表 softmax 变成 top-k 选择加小集合 softmax。

\begin{lstlisting}[numbers=none,caption={top-k 采样流程}]
selected = find_top_k(logits, k)
for token not in selected:
  probability[token] = 0
softmax over selected logits
sample one token from selected distribution
\end{lstlisting}

\code{k=1} 等价于贪心。较大的 \code{k} 给更多多样性。top-k 的工程关键是选择算法。完整排序是 \(\complexity{O(V\log V)}\)，只取前 k 可以用 heap、partial selection 或后端专用 kernel。对每 token decode 来说，词表扫描成本会累积。

\subsection{top-p：按累计概率截断}

top-p 又叫 nucleus sampling。它先按概率从大到小排序，选择最小集合，使累计概率超过 \code{p}：

\[
\sum_{i \in S} p_i \ge p
\]

然后只在集合 \code{S} 中重新归一化采样。top-p 的直觉是：当模型很确定时，候选集合小；当模型不确定时，候选集合大。相比固定 top-k，它能根据分布形态自适应。

代价是需要概率和排序或近似排序。若实现每步对全词表完整排序，sampler 可能成为短 decode 的显著成本。高性能实现通常会结合 top-k 上限、分块选择或 GPU sampler。

\subsection{重复惩罚和频率惩罚}

语言模型容易重复。重复惩罚会根据已生成 token 调整 logits。最简单的策略是：若 token 已出现，把正 logit 除以 penalty，把负 logit 乘以 penalty。这样已出现 token 的相对分数下降。

\begin{lstlisting}[numbers=none,caption={重复惩罚示意}]
for token in generated_tokens:
  if logits[token] > 0:
    logits[token] /= penalty
  else:
    logits[token] *= penalty
\end{lstlisting}

频率惩罚则按出现次数扣分。存在惩罚只看是否出现过。它们都是采样策略，不是模型本体。业务上要谨慎：代码、法律、医学、术语解释中，重复某些 token 可能是正确行为；过强惩罚会破坏格式、变量名和专有名词。

\subsection{停止条件}

生成循环不能只靠 \code{max\_new\_tokens} 停止。常见停止条件包括：

\begin{itemize}
  \item 生成了模型的 EOS token。
  \item 生成了业务指定 stop token 或 stop sequence。
  \item 达到 \code{max\_new\_tokens}。
  \item 达到 \code{max\_context\_tokens}。
  \item 请求超时或被取消。
  \item sampler 发现非法分布，例如全部概率为零。
\end{itemize}

stop sequence 比 stop token 更难，因为它可能跨多个 token。服务端要在 token 流里维护一个小窗口，检查最近 token 是否匹配 stop sequence。若流式输出已经把 stop sequence 的一部分发给客户端，可能需要延迟输出或在业务协议中接受这种行为。不同产品会选择不同策略，但必须明确。

\subsection{随机数和可复现性}

随机采样必须有可控随机数源。固定 seed、固定 sampler 参数、固定 logits，应该得到同一 token 序列。多线程服务中不要使用全局共享随机数直接采样，否则请求之间会互相影响。

\begin{lstlisting}[numbers=none,caption={SamplerState 的最小字段}]
SamplerState:
  seed
  step
  temperature
  top_k
  top_p
  repetition_penalty
  generated_token_counts
\end{lstlisting}

把 \code{step} 放进状态是为了让 trace 可重放。若第 37 步出现异常，可以用同一 logits、同一 sampler state 复现该步选择。业务上，完全可复现有助于审计和回归；创意任务可以使用随机 seed，但也应该把 seed 记录到 trace。

\subsection{beam search：搜索多个候选序列}

beam search 不是每步只保留一个 token，而是保留多个候选序列。每个候选有累计 log 概率。每一步扩展所有候选，保留得分最高的 \code{beam\_width} 条。

\begin{lstlisting}[numbers=none,caption={beam search 伪代码}]
beams = [empty sequence with score 0]
for step in 0..max_new_tokens:
  candidates = []
  for beam in beams:
    logits = run_model(beam.sequence)
    for token in top_candidates(logits):
      candidates.add(beam + token, updated_score)
  beams = best beam_width candidates
\end{lstlisting}

beam search 常用于翻译、较短结构化任务或需要高概率序列的场景。但对聊天式大模型，它可能输出更保守、更模板化的文本，而且计算成本明显更高。若没有共享 KV cache 和批量扩展，beam width 为 \code{B} 就接近把 decode 成本放大 \code{B} 倍。实现 beam search 时，每个 beam 都有自己的 token 序列和 KV cache 分支，不能简单共享一个 session 状态。

\subsection{投机解码：用小模型减少大模型步数}

投机解码的目标是减少大模型 decode 次数。思路是用一个小模型先快速提出多个候选 token，再让大模型一次验证这些候选。若候选被接受，就相当于大模型一次前进多个 token；若被拒绝，则回退到大模型给出的分布。

简化流程：

\begin{lstlisting}[numbers=none,caption={投机解码简化流程}]
draft_tokens = small_model.generate_k_tokens(context)
big_model evaluates context + draft_tokens
for each draft token:
  accept if verification rule passes
  otherwise sample from corrected distribution and stop draft chain
\end{lstlisting}

投机解码不是简单“用小模型代替大模型”。它的关键是验证规则要保持目标大模型的分布不变或可控近似。工程上还要处理两套模型的 tokenizer 一致性、KV cache、显存/内存占用、调度和 trace。它适合在 decode 被大模型单步开销限制时使用，但会增加系统复杂度。

\subsection{采样器放 CPU 还是 GPU}

如果 logits 在 GPU 上，采样器有两种选择。把 logits 拷回 CPU 采样，实现简单，但每步有设备到主机拷贝和同步。把 sampler 放在 GPU，能减少同步，但实现 top-k/top-p、随机数和 stop 检查更复杂。

\begin{center}
\begin{tabular}{p{0.20\linewidth}p{0.34\linewidth}p{0.30\linewidth}}
\toprule
位置 & 优点 & 风险 \\
\midrule
CPU sampler & 简单、易调试、业务逻辑方便 & 每步 copy/sync，GPU decode 受阻 \\
GPU sampler & 减少同步，适合高吞吐 & 实现复杂，trace 和可复现性更难 \\
混合策略 & top-k 在 GPU，最终选择在 CPU & 仍需管理拷贝和候选集合 \\
\bottomrule
\end{tabular}
\end{center}

第一版可以先用 CPU sampler，但 trace 必须记录 logits copy 和 sampler 时间。否则接入 GPU 后，会误以为 kernel 很快就代表端到端快。

\subsection{采样错误怎样定位}

若生成结果异常，不要先改温度。按下面顺序定位：

\begin{enumerate}
  \item 固定 sampler 为贪心，检查 logits top-1 是否和 reference 一致。
  \item 固定 seed，检查同一 logits 下 sampler 是否复现。
  \item 检查 temperature、top-k、top-p 的应用顺序。
  \item 检查重复惩罚是否只作用于已生成 token，而不是 prompt 或全部词表。
  \item 检查 stop token 和 stop sequence 是否在 token 边界处理。
  \item 检查 GPU/CPU sampler 是否使用同一随机数和同一排序规则。
\end{enumerate}

采样器看起来是“小组件”，但它决定最终可见输出。工程报告应该把 \code{sampler\_profile} 记录到 trace，否则两次输出不同可能根本不是模型或后端问题。

\subsection{本节验收线}

读完本节，应该能实现一个可控采样器：

\begin{itemize}
  \item 知道 logits、softmax、概率和 token 选择的关系。
  \item 能实现稳定 softmax、贪心、温度、top-k、top-p 和重复惩罚。
  \item 能解释采样策略对确定性、多样性、质量和延迟的影响。
  \item 能设计 \code{SamplerState}，固定 seed 后复现生成序列。
  \item 能区分 beam search 和普通采样的 KV cache 成本。
  \item 能说明投机解码为什么需要验证，而不是直接相信小模型。
  \item 能把 sampler 时间、logits copy、seed 和参数写进 trace。
\end{itemize}
